\documentclass[11pt,a4paper]{article}
\usepackage[english]{fga}

\title{\textbf{Function-generated arithmetics}\\[2pt]
\large Catalogue, structural classification and numerical limits of the scheme
$a \gop b = \gf^{-1}\big(\gf(a)+\gf(b)\big)$}
\author{Maciej Szulia\\[2pt]\normalsize\texttt{maciej.szulia@gmail.com}}
\date{November 2026}

\begin{document}
\maketitle

\begin{abstract}
\noindent
This note studies the family of operations $a\gop b=\gf^{-1}\big(\gf(a)+\gf(b)\big)$,
where $\gf$ is a bijection. Every such operation is a transport of the structure
of the additive group $(\mathbb{R},+)$ along $\gf$. Associativity, commutativity,
the neutral element, additive inverses and the distributivity of the induced
multiplication therefore hold automatically. We give a catalogue of ten
generators with exact closed forms and classify the induced structure as a
group, a monoid or a semigroup. The classification depends on the image of
$\gf$, not on the shape of the formula. For the two generators without an
elementary inverse (Lambert and Kepler) we measure the numerical cost. For the
synergy operator with the term $2\sqrt{sX}$ we prove that for $s>0$ it has no
neutral element and is not associative, so it is not an arithmetic in the sense
of Grossman and Katz; we report the measured size of those violations. We close
with the physical laws that literally have this form, and with a list of claims
that the scheme does not support.

\smallskip
\noindent\textbf{Keywords:} non-Newtonian calculus; transport of structure;
Kolmogorov--Nagumo mean; Lambert $W$ function; numerical verification.
\end{abstract}

\section{Introduction}

Many composition laws in physics share one form. First one adds a conserved
quantity. Then one reads off the measured quantity. The resistance of two
resistors in parallel is not the sum of the resistances. The focal length of two
lenses is not the sum of the focal lengths. The root-mean-square of two voltages
is not the sum of the root-mean-square values. In each case the addition happens
in a different scale.

One scheme describes all of them:
\begin{equation}
\label{eq:base}
a\gop b=\gf^{-1}\big(\gf(a)+\gf(b)\big).
\end{equation}
The function $\gf$ is a bijection. We call it the generator. The operation $\gop$
is addition in the scale $\gf$.

The note answers three questions.
\begin{enumerate}[label=(\roman*),leftmargin=2.2em]
  \item Which part of \eqref{eq:base} is a theorem, and which part is a choice?
  \item When is the induced structure a group, and when only a monoid or a
        semigroup?
  \item What happens when a term describing the environment, the so-called
        synergy term, is added to the sum?
\end{enumerate}

The answers are largely negative. Scheme \eqref{eq:base} does not create a new
arithmetic. It transports an existing one. The synergy term leaves the algebra
of groups. Neither fact discards the scheme. They place it where it is useful:
as a uniform notation and as a ready set of formulas.

\section{The scheme and the transport theorem}

\begin{definition}[function-generated arithmetic]
\label{def:arith}
Let $D\subseteq\mathbb{R}$ and let $\gf\colon D\to\mathbb{R}$ be a bijection.
For $a,b\in D$ we define
\[
a\gop b=\gf^{-1}\big(\gf(a)+\gf(b)\big),\qquad
a\gsub b=\gf^{-1}\big(\gf(a)-\gf(b)\big).
\]
Multiplication by a scalar $n\in\mathbb{R}$ and division by $n\neq0$ are
\[
n\gscale a=\gf^{-1}\big(n\,\gf(a)\big),\qquad
a\gdiv n=\gf^{-1}\!\left(\frac{\gf(a)}{n}\right).
\]
Full multiplication and full division of two elements are
\[
a\gmul b=\gf^{-1}\big(\gf(a)\,\gf(b)\big),\qquad
a\gdiv_{\!f} b=\gf^{-1}\!\left(\frac{\gf(a)}{\gf(b)}\right).
\]
\end{definition}

\begin{theorem}[transport of structure]
\label{thm:transport}
The structure $(D,\gop)$ is an abelian group. The map $\gf$ is an isomorphism
of groups $(D,\gop)\to(\mathbb{R},+)$.
\end{theorem}

\begin{proof}
$\gf$ is a bijection, so it has an inverse. For $a,b,c\in D$:
$\gf\big((a\gop b)\gop c\big)=\gf(a)+\gf(b)+\gf(c)
=\gf\big(a\gop(b\gop c)\big)$. Commutativity follows from the commutativity of
addition. The neutral element is $0_{\gf}=\gf^{-1}(0)$, because
$\gf(a\gop 0_{\gf})=\gf(a)$. The inverse of $a$ is
$\gsub a=\gf^{-1}\big(-\gf(a)\big)$.
\end{proof}

\begin{corollary}
\label{cor:automatic}
The following properties hold for every generator and need no separate proof:
associativity, commutativity, existence of $0_{\gf}$, existence of additive
inverses, distributivity of $\gmul$ over $\gop$, associativity of $\gmul$, the
power formula $a^{\gmul n}=\gf^{-1}\big(\gf(a)^{n}\big)$, and the mean formula
\[
M_{\gf}(a_{1},\dots,a_{n})=\gf^{-1}\!\left(\frac{1}{n}\sum_{k=1}^{n}\gf(a_{k})\right).
\]
\end{corollary}

Corollary \ref{cor:automatic} is known. The mean $M_{\gf}$ is the
quasi-arithmetic mean of Kolmogorov and Nagumo
\cite{kolmogorov1930,nagumo1930}. The family \eqref{eq:base} was described by
Grossman and Katz \cite{grossman1972} as non-Newtonian calculus. The same scheme
is the standard transport of structure in algebra \cite{aczel1966}. The practical
consequence is this. The sentence ``I discovered a new arithmetic'' means, in
this language, ``I chose a new function $\gf$''. The content is the choice of
$\gf$, not the formula \eqref{eq:base}.

\begin{table}[t]
\centering
\small
\caption{What follows from Theorem \ref{thm:transport} and what is a modelling
choice.}
\label{tab:automatic}
\begin{tabularx}{\textwidth}{@{}lX@{}}
\toprule
\textbf{Status} & \textbf{Item} \\
\midrule
theorem & associativity, commutativity, $0_{\gf}$, additive inverses,
distributivity, power and mean formulas \\
choice & the function $\gf$, the domain $D$, the unit and the sign of the
measured quantity \\
separate theorem & existence of a closed form for $\gfi$, numerical stability,
behaviour as $x\to\partial D$ \\
\bottomrule
\end{tabularx}
\end{table}

\section{Classification: group, monoid, semigroup}

Theorem \ref{thm:transport} assumes that $\gf$ maps $D$ onto all of
$\mathbb{R}$. That assumption often fails. Electrical resistance is positive.
Probability lies in $(0,1)$. Absolute temperature is non-negative. The image of
the generator is then a subset of $\mathbb{R}$, and the induced structure is
weaker. The type of structure depends only on the image of $\gf$.

\begin{proposition}
\label{prop:classification}
Let $\gf\colon D\to\mathbb{R}$ be a bijection onto its image $g(D)$.
\begin{enumerate}[label=(\alph*),leftmargin=2.2em]
  \item If $g(D)=\mathbb{R}$, then $(D,\gop)$ is a group.
  \item If $g(D)=[0,\infty)$, then $(D,\gop)$ is a monoid: $0_{\gf}$ exists,
        additive inverses do not.
  \item If $g(D)=(0,\infty)$, then $(D,\gop)$ is a semigroup: no $0_{\gf}$
        exists in $D$.
\end{enumerate}
\end{proposition}

\begin{proof}
Part (a) is Theorem \ref{thm:transport}. For (b), $0\in g(D)$, so $0_{\gf}$
exists. For $a\neq0_{\gf}$ the element $-\gf(a)$ is negative and does not belong
to $g(D)$; the additive inverse does not exist. For (c), $0\notin g(D)$, so
$\gfi(0)$ does not exist in $D$; the operation has no neutral element.
\end{proof}

\begin{table}[t]
\centering
\small
\caption{Catalogue of generators. Each row gives the exact form of $\gf$, the
exact form of $\gfi$, the domain, the image, the type of structure and the
constants $0_{\gf}$ and $1_{\gf}$. The numbers were produced by
\texttt{scripts/verify.py}.}
\label{tab:catalogue}
\setlength{\tabcolsep}{4pt}
\begin{tabularx}{\textwidth}{@{}l l l l l l l l@{}}
\toprule
\textbf{model} & $\gf(x)$ & $\gfi(y)$ & $D$ & $g(D)$ & \textbf{structure}
& $0_{\gf}$ & $1_{\gf}$ \\
\midrule
harmonic & $1/x$ & $1/y$ & $(0,\infty)$ & $(0,\infty)$ & semigroup & -- & -- \\
power $p$ & $x^{p}$ & $y^{1/p}$ & $[0,\infty)$ & $[0,\infty)$ & monoid & $0$ & $1$ \\
cubic & $x^{3}$ & $\sqrt[3]{y}$ & $\mathbb{R}$ & $\mathbb{R}$ & group & $0$ & $1$ \\
geometric & $\ln x$ & $e^{y}$ & $(0,\infty)$ & $\mathbb{R}$ & group & $1$ & $e$ \\
logistic & $\ln\frac{x}{1-x}$ & $\dfrac{1}{1+e^{-y}}$ & $(0,1)$ & $\mathbb{R}$ & group & $0.5$ & $0.7311$ \\
softplus & $\ln(e^{x}-1)$ & $\ln(e^{y}+1)$ & $(0,\infty)$ & $\mathbb{R}$ & group & $0.6931$ & $1.3133$ \\
Lambert & $x+\ln x$ & $W(e^{y})$ & $(0,\infty)$ & $\mathbb{R}$ & group & $0.5671$ & $1$ \\
relativistic & $\operatorname{artanh}\frac{x}{c}$ & $c\tanh y$
& $(-c,c)$ & $\mathbb{R}$ & group & $0$ & $0.7616c$ \\
Kepler & $x+k\sin x$ & numerically & $\mathbb{R}$ & $\mathbb{R}$ & group
& $0$ & $0.5533$ \\
Tsallis & $\ln(1+\varepsilon x)$ & $\dfrac{e^{y}-1}{\varepsilon}$
& $(-\tfrac1\varepsilon,\infty)$ & $\mathbb{R}$ & group & $0$ & $3.4366$ \\
\bottomrule
\end{tabularx}
\end{table}

Table \ref{tab:catalogue} shows that two models used in practice are weaker than
usually assumed. The power arithmetic with $p=2$ describes root-mean-square
values and kinetic energy. It is a monoid. It has no subtraction: $a\gsub b$
exists only for $a\ge b$. The harmonic arithmetic describes equivalent
resistance and focal length. It is a semigroup. It has no zero. The sentence
``add zero resistance'' has no meaning in that arithmetic.

\section{Models that are physical laws}

In this section $\gop$ is not a modelling choice. It is the content of the law.

\paragraph{Parallel resistance and focal length.}
For the generator $\gf(x)=1/x$ we obtain $a\gop b=\dfrac{ab}{a+b}$. This is the
formula for the equivalent resistance of two resistors in parallel and the
formula for the focal length of two thin lenses. Check: $6\gop3=2.000000$ and
$50\gop200=40.000000$.

\paragraph{Root-mean-square and radiation.}
For $\gf(x)=x^{2}$ we obtain $a\gop b=\sqrt{a^{2}+b^{2}}$. This is the
root-mean-square of two uncorrelated waveforms. Example: $230\gop45=234.360833$
volts. For $\gf(x)=x^{4}$ we obtain the Stefan--Boltzmann law in the form of a
sum of fluxes: $800\gop1000=1089.617257$ kelvin. The value $1088.6$ quoted in
some sources is wrong.

\paragraph{Decibels.}
For $\gf(L)=10^{L/10}$ we obtain
$L_{1}\gop L_{2}=10\log_{10}\big(10^{L_{1}/10}+10^{L_{2}/10}\big)$.
Two sources of $80$ dB give $83.010300$ dB, not $160$ dB.

\paragraph{Composition of velocities.}
For $\gf(v)=\operatorname{artanh}(v/c)$ we obtain
\[
v_{1}\gop v_{2}=\frac{v_{1}+v_{2}}{1+v_{1}v_{2}/c^{2}}.
\]
This is Einstein's formula \cite{einstein1905}. For $v_{1}=v_{2}=0.5c$ the result
is $0.8c$. For $v_{1}=v_{2}=0.9c$ the result is $0.994475c$. The result never
exceeds $c$, because $\tanh$ maps into $(-1,1)$. The structure is known as a
gyrovector space \cite{ungar2008}.

\paragraph{Pooling independent evidence.}
For $\gf(p)=\ln\frac{p}{1-p}$ we obtain
\[
p_{1}\gop p_{2}=\frac{p_{1}p_{2}}{p_{1}p_{2}+(1-p_{1})(1-p_{2})}.
\]
This is the update of the odds for independent evidence in favour of an event.
Example: $0.8\gop0.7=0.903226$. Two concordant independent witnesses raise the
confidence, and the result stays in $(0,1)$.

\begin{table}[t]
\centering
\small
\caption{Physical laws of the form \eqref{eq:base} and their numerical
results.}
\label{tab:laws}
\begin{tabularx}{\textwidth}{@{}l l l X@{}}
\toprule
\textbf{law} & \textbf{generator} & \textbf{meaning of $\gop$} & \textbf{example} \\
\midrule
parallel resistance, focal length & $1/x$ & sum of conductances & $6\gop3=2$ \\
root-mean-square & $x^{2}$ & sum of powers & $230\gop45=234.36$ \\
Stefan--Boltzmann law & $x^{4}$ & sum of fluxes & $800\gop1000=1089.62$ \\
sound level & $10^{L/10}$ & sum of intensities & $80\gop80=83.01$ \\
velocity composition & $\operatorname{artanh}(v/c)$ & sum of rapidities & $0.5c\gop0.5c=0.8c$ \\
independent evidence & $\ln\frac{p}{1-p}$ & sum of logits & $0.8\gop0.7=0.9032$ \\
\bottomrule
\end{tabularx}
\end{table}

\section{Two generators without an elementary inverse}

\subsection{The Lambert generator}

Let $\gf(x)=x+\ln x$ for $x>0$. The function is strictly increasing, so it is a
bijection onto $\mathbb{R}$. Its inverse has the closed form
$\gfi(y)=W(e^{y})$, where $W$ is the Lambert function \cite{corless1996}. The
constants are $0_{\gf}=W(1)=\Omega=0.5671432904097838$ and $1_{\gf}=1$.

The model has two regimes. For $a,b\ll1$ the logarithm dominates, so
\[
a\gop b=ab\,(1+a+b)+O\big((ab)^{2}\big).
\]
Measurement: for $a=b=0.01$ the relative deviation from the product is $0.0201$,
that is $a+b$. For $a,b\gg1$ the identity term dominates, so $a\gop b\to a+b$.
Measurement: the relative deviation from the sum falls below $1\%$ for $x>25.1$
and below $0.1\%$ for $x>251$ (Figure \ref{fig:phase}).

\begin{figure}[t]
\centering
\begin{tikzpicture}
\begin{loglogaxis}[
  fgaaxis,
  xlabel={$x$},
  ylabel={result},
  legend pos=south east,
  xmin=1e-5, xmax=1e5, ymin=1e-6, ymax=1e7,
]
\addplot[very thick, blue!70!black] table[x=x,y=add]
  {data/lambert_phase.dat};
\addlegendentry{$0.8\gop x$}
\addplot[thick, black!55, dashed, domain=1:1e5] {0.8+x};
\addlegendentry{$0.8+x$ (for $x>1$)}
\addplot[thick, red!70!black, dotted, mark=none, domain=1e-5:1] {0.8*x};
\addlegendentry{$0.8\,x$ (for $x<1$)}
\end{loglogaxis}
\end{tikzpicture}
\caption{The Lambert arithmetic with $a=0.8$. For $x\ll1$ the operation
coincides with multiplication. For $x\gg1$ it coincides with addition. The
transition is smooth and has no threshold. Data from
\texttt{paper/data/lambert\_phase.dat}.}
\label{fig:phase}
\end{figure}

The inverse $\gfi$ needs care. The closed form $W(e^{y})$ fails numerically. The
expression $e^{y}$ exceeds the double-precision range for $y>709.78$, and $W$ is
then out of reach. The correct solution does not form $e^{y}$. It solves
$x+\ln x=y$ directly, in the logarithmic domain, by Newton iteration started at
$y-\ln y$. Table \ref{tab:lambert} compares the two ways.

\begin{table}[t]
\centering
\small
\caption{Computing $\gfi(y)$ for the Lambert generator. The column ``closed
form'' is $W(e^{y})$ in IEEE-754 double precision. The column ``asymptotic'' is
$y-\ln y+\ln y/y$, the approximation used in the source code.}
\label{tab:lambert}
\begin{tabular}{@{}rrrr@{}}
\toprule
$y$ & closed form $W(e^{y})$ & solve in $\ln$ domain & asymptotic error \\
\midrule
$10$      & $7.9294200950$  & $7.9294200950$  & $2.2\cdot10^{-4}$ \\
$100$     & $95.4414866456$ & $95.4414866456$ & $6.3\cdot10^{-6}$ \\
$500$     & $493.7978737290$& $493.7978737290$& $1.1\cdot10^{-7}$ \\
$700$     & $693.4583088790$& $693.4583088790$& $4.4\cdot10^{-8}$ \\
$710$     & overflow        & $703.4440117120$& $4.2\cdot10^{-8}$ \\
$10^{4}$  & overflow        & $9990.7905809943$ & $3.3\cdot10^{-11}$ \\
$10^{6}$  & overflow        & $999986.1845032576$ & $1.2\cdot10^{-16}$ \\
\bottomrule
\end{tabular}
\end{table}

Two conclusions. First, the asymptotic expansion has an error of
$2.2\cdot10^{-4}$ at $y=10$ and $1.1\cdot10^{-7}$ at $y=500$. It therefore does
not give machine accuracy, contrary to a claim often repeated. Second, the
overflow is not a defect of $W$. It is a consequence of evaluating $e^{y}$ where
it suffices to solve the original equation.

\subsection{The Kepler generator}

Let $\gf(x)=x+k\sin x$ with $k\in(0,1)$. The derivative $\gf'(x)=1+k\cos x$ is
positive, so $\gf$ is a bijection $\mathbb{R}\to\mathbb{R}$. The inverse has no
closed form. We compute it by Newton iteration started at $y$; the number of
iterations to accuracy $10^{-15}$ is between three and five. The function is
known as Kepler's equation in astronomy.

The constants of this model differ from those quoted in secondary sources.
$0_{\gf}=0$, but $1_{\gf}=0.553315638360866$, not $1$. Check:
$\gf(1)=1.715250337086712$, so $\gf(1)\neq1$. Earlier tables that give
$1_{\gf}=1$ are wrong.

\section{The synergy operator: a negative result}

This section studies the modification in which the sum depends on the
environment.

\begin{definition}[the synergy operator]
\label{def:synergy}
Let $\gf\colon D\to\mathbb{R}$ be a bijection and let $s\ge0$ be a parameter. We
define
\begin{equation}
\label{eq:synergy}
a\gop_{s}b=\gf^{-1}\!\left(\Big(\sqrt{\gf(a)+\gf(b)}+\sqrt{s}\Big)^{2}-s\right).
\end{equation}
Equivalently $a\gop_{s}b=\gf^{-1}\big(X+2\sqrt{sX}\big)$ with
$X=\gf(a)+\gf(b)$ and $X\ge0$.
\end{definition}

\begin{table}[t]
\centering
\small
\caption{Measured violations for \eqref{eq:synergy} with the Lambert
generator, $a=1.3$, $b=1.7$, $c=2.3$. Second column: relative difference between
$(a\gop_s b)\gop_s c$ and $a\gop_s(b\gop_s c)$. Third column:
$|a\gop_s 0_{\gf}-a|$.}
\label{tab:synergy}
\begin{tabular}{@{}rrr@{}}
\toprule
$s$ & associativity violation & neutral-element violation \\
\midrule
$0$     & $0$            & $0$ \\
$0.01$  & $1.07\cdot10^{-2}$ & $0.145$ \\
$0.10$  & $2.83\cdot10^{-2}$ & $0.478$ \\
$0.50$  & $4.82\cdot10^{-2}$ & $1.139$ \\
$1.00$  & $5.75\cdot10^{-2}$ & $1.673$ \\
$2.00$  & $6.64\cdot10^{-2}$ & $2.471$ \\
$5.00$  & $7.64\cdot10^{-2}$ & $4.156$ \\
\bottomrule
\end{tabular}
\end{table}

\begin{figure}[t]
\centering
\begin{tikzpicture}
\begin{loglogaxis}[
  fgaaxis,
  xlabel={$s$},
  ylabel={violation},
  legend pos=north west,
  xmin=0.008, xmax=6, ymin=5e-3, ymax=1e1,
]
\addplot[very thick, blue!70!black, mark=*, mark size=2pt]
  table[x=s,y=assoc_defect] {data/synergy_assoc.dat};
\addlegendentry{associativity violation}
\addplot[very thick, red!70!black, mark=square*, mark size=2pt]
  table[x=s,y=identity_defect] {data/synergy_assoc.dat};
\addlegendentry{neutral-element violation}
\end{loglogaxis}
\end{tikzpicture}
\caption{The violations grow like $\sqrt{s}$. The vertical axis is logarithmic.
As $s\to0$ the operator returns to the arithmetic \eqref{eq:base}.}
\label{fig:synergy}
\end{figure}

\begin{proposition}
\label{prop:synergy}
For $s>0$ the operation \eqref{eq:synergy} is commutative, is not associative,
and has no neutral element in $D$.
\end{proposition}

\begin{proof}
Commutativity follows from the symmetry of $X$. For $a=\gf^{-1}(x)$ and
$b=\gf^{-1}(0)$ we obtain $a\gop_{s}b=\gf^{-1}\big(x+2\sqrt{sx}\big)\neq a$ for
$x>0$, so $0_{\gf}$ is not neutral. For associativity a numerical
counterexample suffices: with the Lambert generator, $a=1.3$, $b=1.7$, $c=2.3$
and $s=0.5$ the two bracketings give $11.626886$ and $12.215317$; the relative
difference is $4.8\cdot10^{-2}$.
\end{proof}

The conclusion is this. Operator \eqref{eq:synergy} is not an arithmetic. It is
not an isomorphism of $(\mathbb{R},+)$, it has no neutral element, and it
depends on the order of composition. It can be used as a one-step nonlinear
update in a control loop. It cannot support the claims that it ``changes the
arithmetic'', that it ``violates conservation of energy'', or that ``the sum of
all doses'' is well defined. The last sentence matters most: without
associativity there is no order-independent result for many doses.

One further observation. The term $2\sqrt{sX}$ is clipped to zero when $X<0$.
For the Lambert generator the condition $X<0$ corresponds to $a<0.5671$. The
operator then acts as the identity. This asymmetry is an artefact of the
clipping, not a physical property.

\section{Claims that the scheme does not support}

The list below collects statements that appear in expositions of scheme
\eqref{eq:base} and do not follow from it.

\begin{enumerate}[leftmargin=1.6em]
  \item \textbf{Kepler's new one equals $1$.} False. $1_{\gf}=0.5533$ for
        $k=0.85$, and the value depends on $k$.
  \item \textbf{Zero in the harmonic arithmetic.} It does not exist in
        $(0,\infty)$; $0_{\gf}=\infty$ lies outside the domain.
  \item \textbf{Subtraction in the power arithmetic.} It exists only for
        $a\ge b$; for $a<b$ the result is complex.
  \item \textbf{Removal of the black-hole singularity by a logarithmic
        substitution.} A change of coordinates does not remove the Kretschmann
        invariant, which for the Schwarzschild metric equals
        $48G^{2}M^{2}/(c^{4}r^{6})$ and diverges at $r=0$ in every frame. The
        tortoise coordinate $r^{*}=r+r_{s}\ln|r/r_{s}-1|$ is a standard tool for
        horizon analysis, not a model of the interior.
  \item \textbf{Warp drive as a consequence of the arithmetic.} The Alcubierre
        metric requires negative energy density. That is an energy-cost
        hypothesis, not a result of the scheme.
  \item \textbf{Necessity of five descriptive elements.} The number of elements
        is an editorial choice. No theorem forces it.
  \item \textbf{Novelty of the scheme.} The scheme is from 1972
        \cite{grossman1972}. What can be new is the choice of $\gf$, the domain
        of application and the implementation.
\end{enumerate}

\section{Conclusions}

Scheme \eqref{eq:base} is a language, not a theory. Its value has three parts.
First, it gives one notation for laws that look different. Second, it carries
ready theorems of algebra into a new context without proof. Third, it exposes
limits: monoid, semigroup and missing inverses are information about the model,
not about the subject.

The synergy term is useful as a controller. It is not an arithmetic. Calling it
an arithmetic leads to false statements about energy conservation and about
independence from the order of composition.

Further work has three directions. First: verification of the models on measured
data rather than on chosen numbers. Second: an implementation of a controller
based on the saturation generator, compared with a PID controller on the same
plant. Third: publication of code and data together with the text, so that every
number in the tables can be reproduced by one command.

{\footnotesize
\setlength{\itemsep}{1pt}
\begin{thebibliography}{9}
\bibitem{grossman1972} M. Grossman, R. Katz.
\emph{Non-Newtonian Calculus}. Lee Press, Pigeon Cove, 1972.
\bibitem{kolmogorov1930} A. N. Kolmogorov.
Sur la notion de la moyenne. \emph{Atti Accad. Naz. Lincei} 12 (1930), 388--391.
\bibitem{nagumo1930} M. Nagumo.
Über eine Klasse von Mittelbildungen mit Anwendungen auf die
Determinantentheorie. \emph{Tôhoku Math. J.} 31 (1930), 71--79.
\bibitem{aczel1966} J. Aczél.
\emph{Lectures on Functional Equations and Their Applications}.
Academic Press, New York, 1966.
\bibitem{einstein1905} A. Einstein.
Zur Elektrodynamik bewegter Körper. \emph{Ann. Phys.} 17 (1905), 891--921.
\bibitem{ungar2008} A. A. Ungar.
\emph{Analytic Hyperbolic Geometry and Albert Einstein's Special Theory of
Relativity}. World Scientific, 2008.
\bibitem{tsallis1988} C. Tsallis.
Possible generalization of Boltzmann--Gibbs statistics.
\emph{J. Stat. Phys.} 52 (1988), 479--487.
\bibitem{corless1996} R. M. Corless, G. H. Gonnet, D. E. G. Hare, D. J. Jeffrey,
D. E. Knuth. On the Lambert $W$ function.
\emph{Adv. Comput. Math.} 5 (1996), 329--359.
\bibitem{czachor2016} M. Czachor.
Relativity of arithmetic as a fundamental symmetry of physics.
\emph{Quantum Stud. Math. Found.} 3 (2016), 123--133.
\end{thebibliography}
}

\end{document}
